\documentclass[11pt,a4paper]{article}
\usepackage{iftex}
\ifPDFTeX
\usepackage[utf8]{inputenc}
\usepackage[T1]{fontenc}
\usepackage{lmodern}
\else
\usepackage{fontspec}
\setmainfont{lmroman10-regular.otf}[BoldFont=lmroman10-bold.otf,ItalicFont=lmroman10-italic.otf,BoldItalicFont=lmroman10-bolditalic.otf]
\setmonofont{lmmono10-regular.otf}
\fi
\usepackage[ngerman]{babel}
\usepackage[a4paper,margin=22mm]{geometry}
\usepackage{microtype}
\usepackage{amsmath,amssymb}
\usepackage{booktabs,longtable,tabularx,array}
\usepackage{enumitem}
\usepackage{xcolor}
\usepackage{xurl}
\usepackage{hyperref}

\definecolor{navy}{HTML}{12334A}
\definecolor{teal}{HTML}{146B75}
\definecolor{softgrey}{HTML}{F1F3F5}
\hypersetup{colorlinks=true,linkcolor=navy,citecolor=teal,urlcolor=teal,
  pdftitle={Warum die Daten Fehler enthalten -- und wie Regeln, KI und Mensch sie abfangen},
  pdfauthor={ETHack 2026}}
\urlstyle{same}
\newcolumntype{P}[1]{>{\raggedright\arraybackslash}p{#1}}
\newcommand{\tco}{\ensuremath{\mathrm{t\,CO}_{2}}}
\setlength{\parindent}{0pt}
\setlength{\parskip}{5pt}
\setlength{\tabcolsep}{4pt}
\setlength{\emergencystretch}{2em}
\setlist{itemsep=3pt,topsep=4pt,leftmargin=*}
\renewcommand{\arraystretch}{1.15}
\setcounter{tocdepth}{2}

\newcommand{\befund}[1]{%
  \par\vspace{2pt}\noindent\colorbox{softgrey}{%
  \begin{minipage}{\dimexpr\linewidth-2\fboxsep}\small\vspace{3pt}#1\vspace{3pt}\end{minipage}}\par\vspace{4pt}}

\newcommand{\statFlags}{382}
\newcommand{\statFehler}{239}
\newcommand{\statPruefen}{143}
\newcommand{\statFirmenBetr}{176}
\newcommand{\statRankRel}{106}
\newcommand{\kiModell}{gpt-5.4-mini}
\newcommand{\kiEskModell}{gpt-5.5}
\newcommand{\kiBestaetigt}{239}
\newcommand{\kiWiderspricht}{0}
\newcommand{\kiAuto}{136}
\newcommand{\kiMensch}{246}
\newcommand{\kiEskaliert}{23}
\newcommand{\kiTokens}{301\,086}
\newcommand{\kiUrteilFehler}{358}
\newcommand{\kiUrteilPlausibel}{16}
\newcommand{\kiUrteilUnklar}{8}
\newcommand{\refN}{10}
\newcommand{\refUrteil}{10}
\newcommand{\refUrsache}{7}
\newcommand{\vEinsWiderspruch}{38}
\newcommand{\vEinsAutoFrei}{33}
\newcommand{\vZweiWiderspruch}{0}
\newcommand{\vZweiAutoFrei}{0}
\newcommand{\vEinsRefUrteil}{9}
\newcommand{\vEinsRefUrsache}{7}
\newcommand{\vEinsFehler}{267}
\newcommand{\vEinsPlausibel}{81}
\newcommand{\vEinsUnklar}{34}


\title{\textcolor{navy}{Warum die Daten Fehler enthalten}\\[4mm]
\large Ursachen, Belege und ein Prüfablauf aus Regeln, KI und Mensch}
\author{ETHack 2026 -- Folgebericht zum Verbesserungsbericht}
\date{12.~September 2026}

\begin{document}
\maketitle

\begin{abstract}
\noindent
Der Verbesserungsbericht hat 380 auffällige Werte beschrieben. Dieser Bericht fragt, \emph{warum} sie entstehen, prüft jede Vermutung an den Rohdaten und an der Dokumentation der Quellen, und baut daraus einen Prüfablauf. Das Ergebnis ist unbequem und hilfreich zugleich: Die meisten Fehler sind keine Tippfehler, sondern entstehen an wenigen Stellen, an denen die Pipeline eine Quelle anders liest, als sie gemeint ist -- eine Null, die „nicht gemeldet“ heißt; ein Messprogramm, das CO\textsubscript{2} gar nicht erfasst; eine Anlage, die beim Zusammenführen doppelt zählt; eine Tochterfirma, die zum falschen Jahr gehört. Zwei Diagnosen des Verbesserungsberichts halten der Prüfung nicht stand: Die CAMPD-Nullen sind keine Fehlzuordnungen, und die ECHO-Werte über 12 sind keine falsch gezählten Programme, sondern doppelt verbundene Zeilen. Ein neues Regelwerk findet \statFlags{} Kandidaten. Ein Sprachmodell beurteilt jeden mit einem Belegpaket; an zehn Fällen mit belegter Antwort liegt es \refUrteil{} von \refN{} Mal richtig. Wer am Ende entscheidet, legt eine feste Regel fest und nicht das Modell.
\end{abstract}

\tableofcontents
\newpage

% ---------------------------------------------------------------------------
\section{Was geprüft wurde}

Grundlage sind der gesammelte Datensatz (\texttt{dataset\_long.parquet}), die Firmen-Zwischentabellen und -- anders als beim Verbesserungsbericht -- die Rohdaten aller Quellen im Zwischenspeicher. Jede Ursache in Abschnitt~\ref{sec:ursachen} ist an mindestens einem konkreten Fall in den Rohdaten nachgesehen und, wo es um die Bedeutung eines Feldes geht, an der Dokumentation der Quelle belegt.

\begin{center}\footnotesize
\begin{longtable}{l P{5.2cm} c r r P{4.2cm}}
\toprule
\textbf{Regel} & \textbf{Befund} & \textbf{Stufe} & \textbf{Fälle} & \textbf{Firmen} & \textbf{Beispiele} \\
\midrule
\endfirsthead
\toprule
\textbf{Regel} & \textbf{Befund} & \textbf{Stufe} & \textbf{Fälle} & \textbf{Firmen} & \textbf{Beispiele} \\
\midrule
\endhead
\texttt{F01} & Zwei Ticker teilen eine CIK: Konzernwerte doppelt & 3 & 0 & CIK-Deduplizierung in consolidate.py \\
\texttt{F02} & Exakt 0 t trotz zugeordneter Anlage & 83 & 0 & fehlend statt 0 in canonical.py \\
\texttt{F03} & 0 t CO2, Programm verlangt kein CO2 & 99 & 0 & CO2 nur aus ARP, RGGI, NSPS4T \\
\texttt{F04} & Mehr als 12 Verstoßquartale je Anlage & 2 & 0 & Join entdoppelt, nur V und S gezählt \\
\texttt{F05} & Mehr als 1,3 t CO2 je MWh & 1 & 0 & Grenze 1,3 t/MWh beim Export \\
\texttt{F06} & Stromrate eines Nicht-Versorgers & 23 & 2 & Rate nur für Kraftwerksflotten \\
\texttt{F07} & Trend oder Basisjahr durch Anlagenbasis verzerrt & 7 & 0 & Trend erst ab drei Jahren mit stabiler Basis \\
\texttt{F08} & Tochter vor oder ohne Konzernzugehörigkeit & 20 & 0 & Gültigkeitszeiträume in resolve.py \\
\texttt{F09} & Unfallrate auf unplausiblen Stunden & 4 & 0 & Stundennenner 200-4.000 je Beschäftigtem \\
\texttt{F10} & Lohnverfahren nur über den Handelsnamen & 0 & 7 & bleibt: Franchise ist ein Grenzfall \\
\texttt{F10} & Lohnverfahren: Team-Zuordnung nicht nachvollziehbar & 35 & 0 & eigene Zuordnung statt Team-Datei \\
\texttt{F11} & Lohnverfahren ohne Nachzahlung und Betroffene & 69 & 0 & Verstöße und Bußgelder als eigene Kennzahlen \\
\texttt{F12} & SBTi: validiert und zurückgezogen vermischt & 25 & 0 & Status je Zieltyp getrennt \\
\texttt{F13} & SBTi: Zieljahr vorbei & 4 & 0 & abgelaufenes Ziel als eigene Kennzahl \\
\texttt{P01} & Extrem im Branchenvergleich & 1 & 1 & bleibt: echte Extreme sind kein Fehler \\
\texttt{P02} & Rangband breiter als 70 Punkte & 6 & 0 & ab Band 70 kein Median mehr \\
\bottomrule%
\end{longtable}
\end{center}

\emph{Fehler} heißt: nachweislich falsch oder irreführend, nicht verrechnen. \emph{Prüfen} heißt: auffällig, aber ohne weiteren Beleg nicht entscheidbar. \statFirmenBetr{} Firmen sind betroffen, \statRankRel{} Fälle betreffen Firmen im Umwelt-Ranking.

% ---------------------------------------------------------------------------
\section{Warum die Fehler entstehen}
\label{sec:ursachen}

Sieben Mechanismen erklären fast alles. Keiner davon ist ein Tippfehler in einer Quelle; fast alle entstehen an der Nahtstelle zwischen dem, was eine Quelle meint, und dem, was die Pipeline daraus liest.

\subsection{Eine Null, die „nicht gemeldet“ heißt}

\textbf{Befund.} 83 Firmenjahre stehen bei exakt 0\,\tco{}, obwohl der Firma Anlagen zugeordnet sind. Kraft Heinz fällt von 45\,412~t (2021) auf 0 (2022), mit denselben vier Anlagen.

\textbf{Mechanismus.} Die EPA-Tabelle führt eine Anlage weiter in den Stammdaten, auch wenn sie keine Menge mehr meldet. Die Pipeline summiert fehlende Mengen als 0. In den Rohdaten sind es 19\,155 Anlagenjahre mit 0~t; bei 17\,394 davon gibt es überhaupt keine gemeldete Menge -- weder direkt noch geliefert -- und bei 15\,733 fehlt sogar der Anlagentyp. Eine regulatorische Erklärung liefert die Aussteigeregel 40~CFR~98.2(i): Wer fünf Jahre unter 25\,000~t oder drei Jahre unter 15\,000~t liegt, darf die Meldung einstellen \cite{ecfr982}. Die Anlage verschwindet dann nicht aus dem Register, ihre Menge schon.

\befund{\textbf{Folge.} Eine Null in dieser Quelle bedeutet „nicht gemessen“, fast nie „emissionsfrei“. Wer sie als Zahl verrechnet, macht Firmen sauberer und erzeugt Trends, die nur das Meldeverhalten abbilden.}

\subsection{Ein Programm, das die Größe nicht misst}

\textbf{Befund.} Marathon Petroleum, Valero, ADM und Merck stehen acht Jahre lang bei 0\,\tco{} gemessenen Kraftwerks-CO\textsubscript{2}. Der Verbesserungsbericht vermutete Fehlzuordnungen.

\textbf{Mechanismus.} Die Zuordnung stimmt -- es sind echte Kessel und Turbinen dieser Firmen. Laut Programmfeld melden sie aber ausschließlich im Programm SIPNOX. Der Datenleitfaden der EPA nennt für SIPNOX „Heat input and NO\textsubscript{X} emissions“, für RGGI dagegen „Heat input and CO\textsubscript{2} emissions“; CO\textsubscript{2} verlangen das Acid Rain Program, RGGI und NSPS4T \cite{camdguide}. Merck meldet 2{,}2 Millionen MWh Last bei 0~t CO\textsubscript{2} -- die Last wird gemessen, das CO\textsubscript{2} nicht.

\befund{\textbf{Korrektur am Verbesserungsbericht.} Keine Fehlzuordnung, sondern eine fehlende Messpflicht. Die richtige Reparatur ist nicht, die Anlagen zu verwerfen, sondern CO\textsubscript{2}-Werte nur aus Programmen zu übernehmen, die CO\textsubscript{2} verlangen.}

\subsection{Zusammenführen, das Zeilen vervielfacht}

\textbf{Befund.} Tesla: 24 Verstoßquartale je Anlage, Northrop Grumman 13{,}5. In drei Jahren sind höchstens 12 möglich.

\textbf{Mechanismus.} Die ECHO-Historie ist eine Zeichenkette mit genau 12 Zeichen, eines je Quartal: \texttt{\_} beanstandungsfrei, \texttt{V} Verstoß, \texttt{S} schwerer Verstoß, \texttt{U} ungeklärt, Leerzeichen ohne Daten. Tesla hat eine einzige Anlage mit \texttt{SSSSSSSSSSSS}. Beim Verbinden mit den EPA-Anlagen über die FRS-Kennung entstehen zwei Zeilen, weil die EPA für dieselbe Anlage zwei Schreibweisen des Konzernnamens führt. Aus 12 werden 24. Zusätzlich zählt die Pipeline \texttt{U} als Verstoß.

\befund{\textbf{Korrektur am Verbesserungsbericht.} Nicht Programme statt Quartale werden gezählt, sondern dieselbe Anlage zweimal. Die Reparatur ist ein Kardinalitätstest beim Verbinden: Eine Anlage darf nach dem Join genau einmal vorkommen.}

\subsection{Eine Quelle, die ihre Angaben nicht prüft}

\textbf{Befund.} McDonald's hat eine Unfallrate von 21{,}2 -- das Dreifache des Branchen-p99. Grundlage sind vier Meldungen mit im Median 50 Arbeitsstunden je Beschäftigtem im Jahr.

\textbf{Mechanismus.} Die Rate ist Fälle je Stunden. Wer die Stunden falsch einträgt, erzeugt jede beliebige Rate. OSHA schreibt dazu selbst: „OSHA also does not validate the counts of workers, hours, or injury and illness counts“ und „concluding that establishments are the most dangerous or the least dangerous solely based on whether they have the highest or lowest rates from these data would be inappropriate“ \cite{ositaguide}. Über alle Meldungen melden 9\,749 Betriebe mehr als 8\,760 Stunden je Beschäftigtem -- mehr, als ein Jahr Stunden hat -- und 18\,774 weniger als 200.

\subsection{Namen statt Kennungen}

\textbf{Befund.} Die Team-Datei nennt für Air Products 212 Lohnverfahren. In den Rohdaten des Arbeitsministeriums gibt es für Air Products keinen einzigen Fall -- aber 15 Fälle von „APDC Cleaning Services“.

\textbf{Mechanismus.} Die Lohnverfahren tragen keine Firmenkennung, nur Rechts- und Handelsnamen, oft mit Filialzusatz. Ein Abgleich über Namensbestandteile findet das Börsenkürzel „APD“ in einem fremden Firmennamen. Dasselbe Muster zeigte sich bei der eigenen Zuordnung: US Steel wurde zu Steel Dynamics, Talen Energy zu Vistra (10{,}5~Mt 2023, noch im Bestand). Die Forschung zu Entity Matching kommt zu einem ähnlichen Bild: Regelbasierte Verfahren neigen zum Übermatching; auf 755\,540 von Fachleuten gelabelten Paaren erreichen Regeln 91{,}3\,\% F1, ein großes Sprachmodell 99{,}0\,\% \cite{opensanctions}.

\subsection{Zeit, die nicht mitgeführt wird}

\textbf{Befund.} PPL wächst angeblich um 240\,\% im Jahr. Constellation hat 2019 bis 2021 exakt 0~t. WestRock erscheint 2023 unter Smurfit Westrock, Pioneer unter ExxonMobil.

\textbf{Mechanismus.} Die Indexliste kennt nur den heutigen Stand, die Tochtertabelle kennt keine Gültigkeitszeiträume. PPL hat 2018 genau eine zugeordnete Anlage mit 5~kt, ab 2019 neun Anlagen mit 28~Mt -- die Zuordnung springt, nicht die Emission. Konzernumbauten (Exelon-Abspaltung 2022, WestRock-Fusion Juli 2024, Pioneer-Übernahme Mai 2024) werden rückwirkend auf alle Jahre projiziert.

\subsection{Definitionen, die zusammengeworfen werden}

\textbf{Befund.} 25 Indexfirmen sind bei SBTi „validiert“ und „Zusage zurückgezogen“ zugleich. GE hat 4{,}44\,\tco{} je MWh. Alphabet, Fox und News Corp stehen je zweimal im Datensatz.

\textbf{Mechanismus.} SBTi führt Status je Zieltyp; „Commitment removed“ heißt, dass ein Ziel nicht binnen 24 Monaten eingereicht wurde \cite{sbtistatus}. Ein gesetztes Nahziel und eine zurückgezogene Netto-null-Zusage sind kein Widerspruch -- die Kennzahl wirft beides zusammen. Die Stromrate eines Heizkraftwerks ist keine Aussage über den Konzern. Und Aktiengattungen mit gemeinsamer CIK sind ein Emittent, nicht zwei.

\begin{center}\footnotesize
\begin{tabularx}{\textwidth}{P{3.2cm} P{3.6cm} X}
\toprule
\textbf{Mechanismus} & \textbf{Wo es auftritt} & \textbf{Gegenmittel} \\
\midrule
Null heißt nicht gemeldet & EPA GHGRP & fehlend statt 0; Off-Ramp als eigener Status \\
Programm misst Größe nicht & CAMPD & CO\textsubscript{2} nur aus ARP, RGGI, NSPS4T \\
Join vervielfacht & ECHO $\leftrightarrow$ GHGRP & Kardinalitätstest nach jedem Join \\
Quelle prüft nicht & OSHA ITA & Nenner-Plausibilität je Meldung vor der Aggregation \\
Namen statt Kennungen & WHD, GHGRP, CAMPD & CIK, LEI (GLEIF, CC0), EIA-860; KI-Gutachten für Grenzfälle \\
Zeit nicht mitgeführt & Tochtertabelle, Indexliste & Gültigkeit von--bis; historische Indexliste \\
Definitionen vermischt & SBTi, eGRID, CIK & Status je Zieltyp; Kennzahl nur für passende Branche; CIK-Deduplizierung \\
\bottomrule
\end{tabularx}
\end{center}

% ---------------------------------------------------------------------------
\section{Wie man es verbessert: zuerst ohne KI}

KI ist nicht der erste Hebel. Die meisten Fehler oben lassen sich mit Regeln verhindern, die schneller, billiger und reproduzierbar sind. Die Reihenfolge folgt der verbreiteten Praxis von Datenverträgen: Prüfungen an jeder Stelle, an der Daten ihre Form ändern, auffällige Zeilen in Quarantäne statt stiller Löschung \cite{pandera,contracts}.

\begin{enumerate}
\item \textbf{Fehlend ist nicht Null.} Wo eine Quelle keine Menge meldet, bleibt der Wert leer. Das betrifft die Summierung der EPA-Mengen und jede Aggregation danach.
\item \textbf{Messpflicht als Datenfeld.} Das Programmfeld von CAMPD mitführen und CO\textsubscript{2} nur übernehmen, wo ein Programm es verlangt.
\item \textbf{Kardinalitätstests.} Nach jedem Join prüfen, dass Schlüssel nicht vervielfacht werden. Ein einziges Assert hätte die ECHO-Werte über 12 verhindert.
\item \textbf{Nenner vor dem Zähler prüfen.} Unfallraten nur aus Meldungen mit 500 bis 3\,500 Stunden je Beschäftigtem; unplausible Meldungen in Quarantäne.
\item \textbf{Kennungen vor Namen.} CIK für SEC, FRS für EPA, ISIN für SBTi, LEI mit den Konzernbeziehungen von GLEIF -- diese stehen unter CC0 \cite{gleif} und beantworten „wem gehört wer“ ohne Namensabgleich.
\item \textbf{Gültigkeitszeiträume.} Jede Zuordnung in der Tochtertabelle bekommt \texttt{gueltig\_von} und \texttt{gueltig\_bis}; die Indexliste wird historisch geführt.
\item \textbf{Harte Grenzen als Export-Asserts.} ECHO $\leq 12$, eGRID $\leq 1{,}3$\,\tco{}/MWh, keine CIK doppelt, $|$Trend$| < 1$. Der Lauf bricht bei Verletzung ab.
\item \textbf{Flags reisen mit.} Ein beanstandeter Wert wird nicht gelöscht, sondern trägt \texttt{quality\_status}, Regel und Begründung bis in die Oberfläche. Das ist inzwischen umgesetzt.
\end{enumerate}

% ---------------------------------------------------------------------------
\section{Wie man KI wirksam einbaut}

\subsection{Wo KI hilft -- und wo nicht}

Regeln finden Kandidaten zuverlässig, können aber den Einzelfall nicht erklären. Genau dort liegt die Stärke eines Sprachmodells: aus einem Bündel von Belegen die wahrscheinlichste Ursache benennen. Ist „APDC Cleaning Services“ Air Products? Ist eine Stromrate bei Berkshire Hathaway ein Fehler oder echte Versorgertätigkeit? Ist ein Sprung von 5~kt auf 28~Mt Klimaschutz, ein Zukauf oder eine Lücke in der Zuordnung?

Nicht geeignet ist ein Modell dafür, fehlende Zahlen zu liefern, Konzernbeziehungen aus dem Gedächtnis zu behaupten oder allein über Werte zu entscheiden, von denen ein Ranking abhängt.

\subsection{Der Aufbau}

\begin{enumerate}
\item \textbf{Regeln setzen die Schwere.} Das Modell kann einer Regel widersprechen, aber eine Regel-Stufe nicht selbst ändern.
\item \textbf{Belegpaket statt Frage.} Jeder Fall geht mit Wert, Zeitreihe, Branchenvergleich und Rohdaten-Auszug an das Modell. Die Anweisung verbietet, Firmenzahlen aus eigenem Wissen zu ergänzen, und verlangt, allgemeines Wissen zu kennzeichnen.
\item \textbf{Feste Antwortform.} Ein JSON-Schema erzwingt Urteil (fehler, plausibel, unklar), eine Ursache aus einer festen Liste, Handlungsvorschlag, Konfidenz und die Frage, ob ein Mensch nötig ist. „Unklar“ ist ausdrücklich eine gute Antwort.
\item \textbf{Weiterleitung durch eine Regel, nicht durch das Modell.} Unsichere Fälle (Konfidenz unter 0{,}6 oder „unklar“) gehen an ein stärkeres Modell; bleiben sie unsicher, an einen Menschen. Korrekturen mit großer Wirkung -- ab einer Million Tonnen oder mit Einfluss auf das Ranking -- gehen immer an einen Menschen.
\item \textbf{Zwischenspeicher.} Gleicher Fall, gleiche Anweisung, gleiches Modell ergibt dieselbe gespeicherte Antwort. Das macht den Lauf reproduzierbar und einen zweiten Lauf kostenlos.
\item \textbf{Messung an belegten Fällen.} Ein Modellurteil ist nur so viel wert wie seine gemessene Trefferquote. Deshalb gibt es eine Referenzmenge mit Fällen, deren Antwort aus den Rohdaten bekannt ist.
\end{enumerate}

\subsection{Warum die Konfidenz des Modells nicht reicht}

Sprachmodelle als Gutachter sind systematisch zu sicher. Eine Untersuchung an 14 Modellen fand etwa für GPT-4o eine Genauigkeit von 49{,}7\,\% bei einem Kalibrierungsfehler von 39{,}3\,\% -- die Antworten häufen sich bei 90 bis 100\,\% Konfidenz, die Trefferquote liegt weit darunter \cite{overconf}. Zwei Gegenmittel sind belegt: mehrere Urteile zusammenführen, was Genauigkeit und Kalibrierung verbessert, und Schwellen nicht setzen, sondern an gelabelten Fällen kalibrieren. Ein aktueller Ansatz legt die Schwelle so fest, dass die Fehlerrate unter den angenommenen Urteilen mit hoher Wahrscheinlichkeit unter einer Vorgabe bleibt, holt bei Unsicherheit zusätzliche Belege und gibt sonst an Menschen ab \cite{judgeabstain}.

\befund{\textbf{Für dieses Projekt heißt das:} Die Schwelle 0{,}6 ist ein Startwert, keine Kalibrierung. Sie wird erst belastbar, wenn die Referenzmenge wächst -- idealerweise aus den Entscheidungen, die Menschen im Prüfablauf ohnehin treffen.}

% ---------------------------------------------------------------------------
\section{Wo der Mensch gebraucht wird}

Ein Mensch ist teuer. Sein Einsatz lohnt sich dort, wo weder Regel noch Modell die Frage aus den vorhandenen Belegen beantworten können, oder wo ein Fehler viel kostet.

\begin{center}\footnotesize
\begin{tabularx}{\textwidth}{P{3.6cm} X P{3.4cm}}
\toprule
\textbf{Grenzfall} & \textbf{Warum Regel und KI nicht reichen} & \textbf{Was der Mensch braucht} \\
\midrule
Stichtag eines Konzernumbaus & Das Modell kennt Übernahmen, aber nicht verlässlich das Vollzugsdatum & Pressemitteilung, 8-K-Meldung \\
Franchise oder eigener Betrieb & Aus dem Handelsnamen nicht entscheidbar & Franchise-Offenlegung (FDD), Adresse \\
Echter Extremwert oder Fehler & Vistra kann real kohlelastig sein -- oder falsch zugeordnet & Anlagenliste, Geschäftsbericht \\
Regel und KI widersprechen sich & Einer von beiden irrt; das Muster ist lehrreich & beide Begründungen nebeneinander \\
Korrektur mit großer Wirkung & ab 1~Mt oder mit Einfluss auf das Ranking & Rohdaten-Auszug, Vorher-Nachher \\
Neue Quellendefinitionen & SBTi führt seit September 2026 neue Status ein & Definitionsdokument der Quelle \\
Lizenzfragen & Rechtliche Einordnung, nicht technische & Nutzungsbedingungen \\
\bottomrule
\end{tabularx}
\end{center}

\textbf{Der Ablauf.} Eine Entscheidung ist eine Zeile in \texttt{data/review/entscheidungen.csv} mit Flag, Entscheidung, Name, Datum und einem Satz Begründung. Beim nächsten Lauf hat sie Vorrang vor Regel und KI und fließt mit Namen in den Datensatz. Die Prüfliste in der Oberfläche sortiert nach Wirkung: rankingrelevant zuerst, dann nach Tonnen, dann nach Unsicherheit des Modells.

\befund{\textbf{Der wichtigste Effekt ist der Rückfluss.} Jede menschliche Entscheidung ist ein neuer Referenzfall für die Messung der KI -- und jede wiederkehrende Entscheidung ein Kandidat für eine neue Regel. So wird der teure Schritt mit der Zeit seltener, ohne dass man ihm blind vertraut.}

% ---------------------------------------------------------------------------
\section{Ergebnisse dieses Laufs}

\begin{center}\small
\begin{tabular}{lr}
\toprule
Kandidaten aus den Regeln & \statFlags \\
\quad davon Stufe Fehler & \statFehler \\
\quad davon Stufe Prüfen & \statPruefen \\
\midrule
Modell (erster Durchgang) & \texttt{\kiModell} \\
Modell (Eskalation) & \texttt{\kiEskModell} \\
KI-Urteil fehler / plausibel / unklar & \kiUrteilFehler{} / \kiUrteilPlausibel{} / \kiUrteilUnklar \\
KI bestätigt Regel-Fehler & \kiBestaetigt \\
KI hält Regel-Fehler für plausibel & \kiWiderspricht \\
eskaliert & \kiEskaliert \\
automatisch entschieden & \kiAuto \\
an Menschen weitergeleitet & \kiMensch \\
Tokens insgesamt & \kiTokens \\
\bottomrule
\end{tabular}
\end{center}

\subsection{Der erste Durchgang hatte einen Fehler -- in der Frage, nicht im Modell}

Im ersten Durchgang lautete die Definition: „plausibel: Der Wert ist ungewöhnlich, aber durch die Belege erklärt“. Das Modell hat sie wörtlich genommen. Bei den Nullen aus SIPNOX-Anlagen schrieb es sinngemäß: 0\,\tco{} ist plausibel, weil das Programm CO\textsubscript{2} nicht misst. Die Ursache stimmt, das Urteil ist trotzdem falsch -- die Null bleibt als Messwert unbrauchbar. So entstanden \vEinsWiderspruch{} Widersprüche zu Regel-Fehlern, und \vEinsAutoFrei{} davon wären ohne Mensch freigegeben worden, weil das Modell sich sicher war.

Zwei Änderungen: Die Frage heißt jetzt ausdrücklich „Darf dieser Wert so als Messwert verrechnet werden?“, mit dem Hinweis, dass eine erklärte Ursache einen Wert nicht gültig macht. Und eine feste Regel: Widerspricht die KI einer Regel, die „Fehler“ sagt, entscheidet immer ein Mensch. Im zweiten Durchgang bleiben \vZweiWiderspruch{} Widersprüche, automatisch freigegeben werden \vZweiAutoFrei.

\begin{center}\small
\begin{tabular}{P{5.6cm}cc}
\toprule
 & \textbf{Durchgang 1} & \textbf{Durchgang 2} \\
\midrule
KI-Urteil fehler / plausibel / unklar & \vEinsFehler{} / \vEinsPlausibel{} / \vEinsUnklar & \kiUrteilFehler{} / \kiUrteilPlausibel{} / \kiUrteilUnklar \\
Widerspruch zu Regel-Fehler & \vEinsWiderspruch & \vZweiWiderspruch \\
davon automatisch freigegeben & \vEinsAutoFrei & \vZweiAutoFrei \\
Referenz: Urteil richtig & \vEinsRefUrteil{} von \refN & \refUrteil{} von \refN \\
Referenz: Ursache richtig & \vEinsRefUrsache{} von \refN & \refUrsache{} von \refN \\
\bottomrule
\end{tabular}
\end{center}

\textbf{Der zweite Durchgang schlägt ins Gegenteil aus.} Die präzisere Frage hat die Freigaben beseitigt, aber das Modell folgt der Regel jetzt fast immer: \kiUrteilFehler{} von \statFlags{} Fällen nennt es Fehler. Darunter sind Fälle, die keine Fehler sind -- ein breites Rangband ist eine Unsicherheit, kein falscher Wert, und ein Lohnverfahren ohne Nachzahlung kann echt sein. Von 81 Fällen, die im ersten Durchgang „plausibel“ hießen, sind 64 jetzt „Fehler“. Die Referenzmenge zeigt das nicht, weil sie neun echte Fehler und nur einen echten Grenzfall enthält. Der Widerspruch zwischen Regel und KI fällt damit als Signal fast weg.

\befund{\textbf{Lehre.} Die hohe Konfidenz des Modells hat den Fehler nicht angezeigt, sondern verdeckt -- genau das Muster aus der Literatur zur Überkonfidenz. Gefunden hat ihn erst der Abgleich mit der Regel. Deshalb gilt: Ein Widerspruch zwischen zwei unabhängigen Prüfern ist das wertvollste Signal für einen Menschen, und Freigaben dürfen nie allein am Modell hängen. Und: Das Belegpaket enthält die Regelmeldung. Ein Gutachter, der das Urteil der Regel sieht, ist nicht mehr unabhängig. Der nächste Schritt ist ein blinder Durchgang ohne Regeltext, dazu eine Referenzmenge mit ebenso vielen Fehlalarmen wie echten Fehlern.}

\subsection{Referenzfälle (Durchgang 2)}

\begin{center}\footnotesize
\begin{longtable}{l P{3.9cm} P{1.2cm} P{2.3cm} P{3.3cm} r P{1.5cm}}
\toprule
\textbf{Firma} & \textbf{Beleg} & \textbf{erwartet} & \textbf{KI-Urteil} & \textbf{KI-Ursache} & \textbf{Konf.} & \textbf{Weg} \\
\midrule
\endfirsthead
\toprule
\textbf{Firma} & \textbf{Beleg} & \textbf{erwartet} & \textbf{KI-Urteil} & \textbf{KI-Ursache} & \textbf{Konf.} & \textbf{Weg} \\
\midrule
\endhead
GOOGL & GOOG und GOOGL teilen CIK 1652044 & fehler & fehler (ja) & doppelzaehlung (ja) & 0.92 & automatisch \\
KHC & 4 Anlagen gemeldet, keine Emissionszeile mit Menge & fehler & fehler (ja) & null\_\allowbreak{}statt\_\allowbreak{}fehlend (ja) & 0.93 & mensch \\
MPC & Programm SIPNOX, CO2 nicht meldepflichtig & fehler & fehler (ja) & programm\_\allowbreak{}misst\_\allowbreak{}groesse\_\allowbreak{}nicht (ja) & 0.93 & automatisch \\
TSLA & 1 Anlage, Historie SSSSSSSSSSSS, 2 Join-Zeilen durch 2 Namensschreibweisen & fehler & fehler (ja) & doppelzaehlung (ja) & 0.98 & mensch \\
BRK-B & Berkshire Hathaway Energy betreibt echte Versorger; nur GICS sagt Financials & plausibel & plausibel (ja) & branchenstruktur\_\allowbreak{}real (ja) & 0.78 & automatisch \\
GE & 4,44 t/MWh, ein Heizkraftwerk, physikalisch unmöglich als Stromrate & fehler & fehler (ja) & programm\_\allowbreak{}misst\_\allowbreak{}groesse\_\allowbreak{}nicht (\textbf{nein}) & 0.89 & mensch \\
PPL & 2018 1 Anlage 5 kt, ab 2019 9 Anlagen 28 Mt & fehler & fehler (ja) & meldefehler\_\allowbreak{}quelle (\textbf{nein}) & 0.87 & mensch \\
VST & Talen Energy ist nicht Teil von Vistra & fehler & fehler (ja) & zuordnung\_\allowbreak{}falsch (ja) & 0.95 & mensch \\
MCD & 4 Meldungen, Median 50 h je Beschäftigtem & fehler & fehler (ja) & definition\_\allowbreak{}oder\_\allowbreak{}einheit (\textbf{nein}) & 0.94 & mensch \\
APD & 0 Rohfälle für Air Products, 15 Fälle 'APDC Cleaning Services' & fehler & fehler (ja) & zuordnung\_\allowbreak{}falsch (ja) & 0.96 & automatisch \\
\bottomrule%
\end{longtable}
\end{center}

Zehn Fälle sind zu wenig, um eine Schwelle zu kalibrieren. Sie reichen, um grobe Fehlanreize zu sehen -- etwa ob das Modell reflexhaft der Regel folgt, ob es „unklar“ nutzt, und ob es die richtige Ursache benennt, nicht nur das richtige Urteil.

% ---------------------------------------------------------------------------
\section{Grenzen und nächste Schritte}

\begin{itemize}
\item Die Regeln sind so gut wie die Mechanismen, die man kennt. Neue Fehlerarten fängt zuerst nur die Statistik -- und die meldet auch echte Extreme.
\item Die Referenzmenge ist klein und von einer Person zusammengestellt. Sie sollte aus menschlichen Entscheidungen wachsen, mit einer zweiten Person bei strittigen Fällen.
\item Das Modell sieht nur die Belege im Paket. Wo ein Rohdaten-Auszug fehlt, kann es nur „unklar“ sagen -- das ist gewollt.
\item Nächste Schritte in der Reihenfolge der Wirkung: fehlend statt Null in der EPA-Summierung, Programmfilter bei CAMPD, Kardinalitätstest bei ECHO, Nenner-Filter bei OSHA, Gültigkeitszeiträume und GLEIF-Beziehungen für die Zuordnung, danach die Kalibrierung der Weiterleitungsschwelle an einer größeren Referenzmenge.
\end{itemize}

% ---------------------------------------------------------------------------
\begin{thebibliography}{99}\small
\bibitem{ecfr982} U.S. Government Publishing Office: 40 CFR 98.2 -- Who must report?, Absatz (i). \url{https://www.ecfr.gov/current/title-40/chapter-I/subchapter-C/part-98/subpart-A/section-98.2}
\bibitem{camdguide} U.S. EPA, Clean Air Markets Division: CAMD's Power Sector Emissions Data Guide, Juli 2022. \url{https://www.epa.gov/system/files/documents/2022-07/CAMD's\%20Power\%20Sector\%20Emissions\%20Data\%20Guide\%20-\%2007182022.pdf}
\bibitem{ositaguide} U.S. OSHA: Injury Tracking Application (ITA) Data Users Guide. \url{https://www.osha.gov/sites/default/files/ITA_data_users_guide.pdf}
\bibitem{opensanctions} OpenSanctions Pairs: Large-Scale Entity Matching with LLMs, arXiv:2603.11051. \url{https://arxiv.org/abs/2603.11051}
\bibitem{sbtistatus} Science Based Targets initiative: What's changing on the SBTi Target Dashboard; Commitment and Target Statuses, Version 2.0, August 2026. \url{https://sciencebasedtargets.org/blog/whats-changing-on-the-sbti-target-dashboard}
\bibitem{pandera} endjin: Creating Quality Gates in the Medallion Architecture with Pandera. \url{https://endjin.com/blog/creating-quality-gates-in-the-medallion-architecture-with-pandera}
\bibitem{contracts} Data Contracts and Data Quality: Enforcing Reliability in Modern Pipelines. \url{https://aicodeinvest.com/data-contracts-data-quality-enforcement-guide/}
\bibitem{gleif} GLEIF: LEI Data Terms of Use (CC0 1.0); Level 2 Data: Who Owns Whom. \url{https://www.gleif.org/en/meta/lei-data-terms-of-use}
\bibitem{overconf} Overconfidence in LLM-as-a-Judge: Diagnosis and Confidence-Driven Solution, arXiv:2508.06225. \url{https://arxiv.org/abs/2508.06225}
\bibitem{judgeabstain} Judge, Retrieve, or Abstain: Uncertainty-Guarded LLM Judging with Provable Risk Guarantees, arXiv:2608.17994. \url{https://arxiv.org/abs/2608.17994}
\end{thebibliography}

\end{document}
